\subsection{Intervalle}
\label{subsec:fundamentals_intervals}

\begin{DefinitionBox}{Intervall}{fundamentals_interval}
  Ein \emph{Intervall} \EnglishTerm{interval} ist eine Teilmenge $I\subseteq\R$
  mit folgender Eigenschaft: Liegen zwei Zahlen $x,z\in I$, dann liegen auch
  alle Zahlen dazwischen in $I$:
  \[
    x,z\in I,\quad x\leq y\leq z
    \quad\Longrightarrow\quad
    y\in I.
  \]
\end{DefinitionBox}

Ein beschränktes Intervall wird durch zwei Randpunkte $a,b\in\R$ mit $a<b$
festgelegt. Runde Klammern schließen einen Randpunkt aus, eckige Klammern
schließen ihn ein. Je nachdem, welche Randpunkte enthalten sind, unterscheidet
man abgeschlossene, offene und halboffene Intervalle.

\begin{center}
  \begin{tikzpicture}[x=1.25cm,y=1cm,font=\small]
  \foreach \y/\interval in {3/{[a,b]},2/{(a,b)},1/{[a,b)},0/{(a,b]}} {
    \node[anchor=east,font=\bfseries] at (-2.6,\y) {$\interval$};
    \draw[->,black!55] (-2.25,\y) -- (3.1,\y);
    \draw[THMBlue,line width=2.2pt] (-1.25,\y) -- (2.0,\y);
    \node[below=3pt] at (-1.25,\y) {$a$};
    \node[below=3pt] at (2.0,\y) {$b$};
  }

  % Abgeschlossenes Intervall [a,b]
  \fill[THMBlue] (-1.25,3) circle (2.3pt);
  \fill[THMBlue] (2.0,3) circle (2.3pt);

  % Offenes Intervall (a,b)
  \fill[white] (-1.25,2) circle (2.8pt);
  \draw[THMBlue,line width=1pt] (-1.25,2) circle (2.8pt);
  \fill[white] (2.0,2) circle (2.8pt);
  \draw[THMBlue,line width=1pt] (2.0,2) circle (2.8pt);

  % Halboffenes Intervall [a,b)
  \fill[THMBlue] (-1.25,1) circle (2.3pt);
  \fill[white] (2.0,1) circle (2.8pt);
  \draw[THMBlue,line width=1pt] (2.0,1) circle (2.8pt);

  % Halboffenes Intervall (a,b]
  \fill[white] (-1.25,0) circle (2.8pt);
  \draw[THMBlue,line width=1pt] (-1.25,0) circle (2.8pt);
  \fill[THMBlue] (2.0,0) circle (2.3pt);
\end{tikzpicture}

\end{center}

\renewcommand{\arraystretch}{1.35}
\begin{tabularx}{\linewidth}{@{}l X l@{}}
  \toprule
  \textbf{Intervalltyp} & \textbf{Mengenschreibweise} & \textbf{Randpunkte} \\
  \midrule
  abgeschlossen
    & $[a,b]=\{x\in\R\mid a\leq x\leq b\}$
    & $a,b$ enthalten \\
  offen
    & $(a,b)=\{x\in\R\mid a<x<b\}$
    & $a,b$ nicht enthalten \\
  links abgeschlossen
    & $[a,b)=\{x\in\R\mid a\leq x<b\}$
    & nur $a$ enthalten \\
  rechts abgeschlossen
    & $(a,b]=\{x\in\R\mid a<x\leq b\}$
    & nur $b$ enthalten \\
  \bottomrule
\end{tabularx}
\renewcommand{\arraystretch}{1}

\begin{RemarkBox}
  Für offene und halboffene Intervalle ist im deutschsprachigen Raum auch
  folgende Schreibweise üblich:
  \[
    (a,b)=]a,b[,
    \qquad
    [a,b)=[a,b[,
    \qquad
    (a,b]=]a,b].
  \]
\end{RemarkBox}

Ist ein Intervall nur auf einer Seite begrenzt, verwendet man $-\infty$ oder
$\infty$:
\begin{align*}
  (-\infty,b]&=\{x\in\R\mid x\leq b\},
  &(-\infty,b)&=\{x\in\R\mid x<b\},\\
  [a,\infty)&=\{x\in\R\mid a\leq x\},
  &(a,\infty)&=\{x\in\R\mid a<x\}.
\end{align*}

\begin{KeyPointBox}{Unendliche Randpunkte}{fundamentals_infinite_interval_endpoints}
  Die Symbole $-\infty$ und $\infty$ sind keine reellen Zahlen und gehören
  daher nie zum Intervall. An einer unendlichen Seite steht immer eine runde
  Klammer.
\end{KeyPointBox}
